\documentclass[11pt,a4paper]{moderncv}
\moderncvstyle{banking}
\moderncvcolor{blue}
\usepackage[scale=0.75]{geometry}
\usepackage[utf8]{inputenc}
\usepackage[T1]{fontenc}
\usepackage{hyperref}

% -----------------------------
% Dados pessoais
% -----------------------------
\name{Matihus}{Molina Larios}
\address{Lima, Peru}{}
\email{matihus.molina@unmsm.edu.pe}
\social[linkedin]{maticus}
\social[github]{Maticus-7}
\extrainfo{\href{http://lattes.cnpq.br/4261863091619974}{Currículo Lattes: 4261863091619974}}

% -----------------------------
\begin{document} 
	\makecvtitle
	
	% -----------------------------
	% Perfil profissional
	% -----------------------------
	\section{Perfil Profissional}
	\cvitem{}{Sou aluno do último ano do Bacharelado em Matemática Pura na Universidad Nacional Mayor de San Marcos (UNMSM). Minha pesquisa atual concentra-se em Sistemas Dinâmicos e Teoria Ergódica, especificamente no estudo de estabilidades topológicas no círculo ($S^1$). Paralelamente, possuo uma sólida base em métodos numéricos, otimização e aprendizado de máquina (Machine Learning), integrando o rigor analítico com ferramentas computacionais avançadas. Tenho como objetivo principal ingressar em um programa de mestrado de excelência no Brasil para desenvolver pesquisas rigorosas em difeomorfismos de Anosov e dinâmica parcialmente hiperbólica.}
	
	% -----------------------------
	% Formação Académica
	% -----------------------------
	\section{Formação Académica}
	\cventry{03/2022 -- 12/2026 (Previsão)}{Bacharelado em Ciências Matemáticas}{Universidad Nacional Mayor de San Marcos (UNMSM)}{Lima, Peru}{}{Média Geral / Coeficiente de Rendimento: 15.324/20}
	
	% -----------------------------
	% Idiomas
	% -----------------------------
	\section{Idiomas}
	\cvitem{\textbf{Espanhol}}{Nativo}
	\cvitem{\textbf{Inglês}}{Intermédio B1 (Diplomado pelo Centro Cultural LIDEMPERÚ. Em preparação para certificação internacional)}
	\cvitem{\textbf{Português}}{Básico A2}
	
	% -----------------------------
	% Projetos de Investigação e Académicos
	% -----------------------------
	\section{Projetos de Investigação}
	
	\cventry{}{Investigação Atual Independente}{Solução numérica para equações singulares do tipo Lane-Emden por ADM e método de volumes finitos}{}{}{
		\begin{itemize}
			\item Extensão do teorema de existência e unicidade para o problema de Lane-Emden generalizado não homogéneo através de um sistema EDO de primeira ordem $3 \times 3$.
			\item Estimação numérica através de métodos de volumes finitos e contraste com Redes Neuronais (PINNs).
			\item \href{https://vii.imms.pe/documentos/26PosterMolinaLarios.pdf}{[Clique aqui para ver o póster oficial (Apresentado no VII EICM)]}
			\item \href{https://maticus-7.github.io/lane-emden/}{[Clique aqui para ver o repositório no GitHub]}
	\end{itemize}}	
	
	\cventry{02/2026 -- 03/2026}{Otimização Numérica e Aprendizagem de Operadores}{SaC$^3$ -- Summer School and Workshop}{}{}{
		\begin{itemize}
			\item \textbf{Orientador:} Leonidas Gkimisis.
			\item Implementei e documentei em Python algoritmos de otimização matemática analisados no workshop, aplicando controlo de versões com Git/GitHub para garantir código reprodutível.
			\item \href{https://github.com/Maticus-7/Data-Science}{[Clique aqui para ver o repositório no GitHub]}
	\end{itemize}}
	
	\cventry{01/2025 -- 03/2025}{Um primeiro encontro com os sistemas dinâmicos caóticos}{PICV-UNMSM}{}{}{
		\begin{itemize}
			\item \textbf{Orientador:} Dr. Jorge L. Crisóstomo.
			\item Estudo da conjugação topológica entre a família quadrática $Q_c$ e o \textit{shift} de Bernoulli, aplicando conceitos de equações diferenciais, análise funcional e topologia.
			\item \href{http://bit.ly/3Hj8NvJ}{[Clique aqui para ver o documento de investigação]}
	\end{itemize}}
	
	% -----------------------------
	% Eventos Académicos e Internacionalização
	% -----------------------------
	\section{Mobilidade Internacional e Eventos Académicos}
	
	\cventry{Setembro 2026}{13th Heidelberg Laureate Forum (HLF)}{Heidelberg}{Alemanha}{}{Selecionado mundialmente como \textit{Young Researcher}. Participação no fórum internacional que reúne os vencedores da Medalha Fields, Prémio Abel e Prémio Turing com jovens investigadores excecionais.}
	
	\cventry{Agosto 2026}{Programa ``San Marquinos pelo Mundo'' (Universidade do Chile)}{Santiago}{Chile}{}{Estágio académico internacional selecionado por ordem de mérito (Top 2 da Escola Profissional). Curso de especialização (40 horas): \textit{Desafios Globais, Transformação Digital e Inovação}.}
	
	\cventry{Dezembro 2025}{VII Encontro Internacional de Ciências Matemáticas (EICM)}{Huaraz}{Peru}{}{Palestrante. Apresentação do póster de investigação sobre soluções numéricas para a equação de Lane-Emden.}
	
	% -----------------------------
	% Competências Técnicas
	% -----------------------------
	\section{Competências Técnicas}
	\cvitem{\textbf{Linguagens e Ambientes}}{Python, R, MATLAB, LEAN4, Git, \LaTeX, Power BI.}
	\cvitem{\textbf{Bibliotecas}}{NumPy, SciPy, Pandas, Scikit-Learn, PyTorch, Matplotlib.}
	
	% ------------------------------
	% Prémios e Distinções
	% ------------------------------
	\section{Prémios e Distinções}
	\cvitem{2025}{Bolsa Permanência de Estudos Nacionais, concedida pelo PRONABEC devido a elevado rendimento académico.}
	\cvitem{2025}{Reconhecimento de Excelência Académica concedido pela Faculdade de Ciências Matemáticas da UNMSM.}  
	
	% -----------------------------
	% Atividades extracurriculares
	% -----------------------------
	\section{Liderança e Atividades Extracurriculares}
	\cvitem{\textbf{Atual}}{Diretor da área de Projetos e Investigação da \textit{IEEE Computational Intelligence Society} (UNMSM) e \textit{Trainee Researcher} do grupo BILDHEIT.}
	\cvitem{\textbf{Anterior}}{Ex-membro da equipa universitária de judo (UNMSM) e do grupo artístico SALSUNI.}
	
	% -----------------------------
	% Certificados
	% -----------------------------
	\section{Certificados e Portefólio} 
	\cvitem{}{ \href{https://maticus-7.github.io/Certificates/}{Os certificados de cursos, perfil Lattes e realizações académicas estão disponíveis no meu repositório digital (Clique aqui)}.}
	
\end{document}